% Deutsche Parallelfassung zu geometry/curvature_space.tex

\chapter{Krümmungsraum}
\label{chap:curvature-space}

\section*{Dasselbe Alignment, intrinsisch gelesen}

Betrachtet wird der identifizierte Gerade--Übergang--Bogen-Zustand aus
Abbildung~\ref{fig:alignment-object-stationing-curvature}. Seine geordneten
Stationsbereiche seien
\[
I_G=[s_0,s_1],\qquad I_U=[s_1,s_2],\qquad I_B=[s_2,s_3].
\]
Dabei ist $s$ die intrinsische Längsposition, weder eine $x$-Koordinate
noch eine dargestellte Kilometerangabe. Die Orientierung legt die
Durchlaufrichtung fest. Ordnung und $s$ bestimmen gemeinsam, welches
Element auf welches folgt, bevor eine geografische Platzierung gewählt wird.

\begin{principle}[Krümmungsraumprinzip]
\label{princ:curvature-space}
Innerhalb von AIM wird die Plangeometrie eines identifizierten
Alignment-Zustands primär durch sein Krümmungsgesetz über der intrinsischen
Längsposition gelesen.
\end{principle}

\section{Krümmung macht die Elementfolge sichtbar}

\begin{definition}[Krümmungsgesetz]
\label{def:curvature-law}
Ein Krümmungsgesetz ist eine Funktion
\[
\kappa:I\to\mathbb R,
\]
die auf einem intrinsischen Bogenlängenintervall
$I\subset\mathbb R$ definiert ist.
\end{definition}

Für das laufende Beispiel gilt
\[
\kappa(s)=
\begin{cases}
0, & s\in I_G,\\
\kappa_U(s), & s\in I_U,\\
\kappa_B, & s\in I_B.
\end{cases}
\label{eq:straight-transition-arc-curvature}
\]
Das Krümmungsband liest die Konstruktion unmittelbar: Null bedeutet Gerade,
konstante von null verschiedene Krümmung bedeutet Kreisbogen, veränderliche
Krümmung bedeutet Übergang. Elementgrenzen erscheinen bei $s_1$ und $s_2$,
wo sich Werte und Ableitungen vergleichen lassen. In der Plandarstellung
können zwei Familien nahezu gleich aussehen und dennoch verschiedene
Endableitungen verbergen.

Das CurvatureBand der App ist eine Implementierungsrealisierung dieser Lesart.
Das mathematische Argument hängt nicht von dieser Oberfläche ab: Jede
Auswertung, die $s$, Elementbereiche, $\kappa$ und die relevanten
Ableitungen erhält, kann dieselbe Struktur offenlegen.

\begin{principle}[Curvature-first-Prinzip]
\label{princ:curvature-first}
Krümmung ist gegenüber rekonstruierten Plankoordinaten primär, bleibt aber ein
Zustand eines identifizierten Alignments statt einer frei schwebenden Kurve.
\end{principle}

\section{Rekonstruktion vor geografischer Platzierung}

\begin{definition}[Intrinsische Rekonstruktion]
\label{def:intrinsic-curvature-reconstruction}
Für $\kappa(s)$ und eine lokale Anfangspose
$(\gamma(s_0),\theta(s_0))$ erfüllen Tangentenrichtung und Plankurve
\[
\theta'(s)=\kappa(s),\qquad
\gamma'(s)=\mathbf t(s)=
\begin{bmatrix}\cos\theta(s)\\\sin\theta(s)\end{bmatrix}.
\]
\end{definition}

Die Integration rekonstruiert eine lokale kartesische Kurve. Eine Änderung
des Ursprungs oder Drehung dieses lokalen Koordinatenrahmens verändert die
Koordinatenrepräsentation, nicht intrinsische Folge oder Krümmungsgesetz. Ein
späterer Realization Context kann denselben Zustand in einem projizierten CRS
oder auf einer Erdreferenzfläche platzieren. Eine physische Realisierung kann
ihn anschließend unvollkommen instanziieren. Diese Abhängigkeiten fügen
metrische, geografische und physische Bedeutung hinzu, ohne die intrinsische
Geometrie des Alignments rückwirkend zu erzeugen.

\section{Krümmungsraum und die erforderlichen Unterscheidungen}

\begin{definition}[Krümmungsraum]
\label{def:curvature-space}
Der Krümmungsraum bezeichnet den Raum der Krümmungsgesetze
\[
\mathcal K=\{\kappa:I\to\mathbb R\},
\]
in dem angegebene Anforderungen an Zulässigkeit, Glattheit, Realisierung und
Betrieb anwendbare Teilmengen auswählen.
\end{definition}

Das Symbol (mathcal K) macht nicht jedes mathematisch ausdrückbare Gesetz
zulässig, anwendbar oder autoritativ. Es stellt den Vergleichsraum bereit, in
dem konstante, veränderliche, beschränkte, differenzierbare und hybride Gesetze
für den identifizierten Ingenieurzustand unterschieden werden.

\begin{definition}[Konstante Krümmung]
\label{def:constant-curvature}
Ein konstantes Krümmungsgesetz erfüllt $\kappa(s)=\kappa_0$. Der Fall
$\kappa_0=0$ ist eine Gerade; $\kappa_0\neq0$ rekonstruiert einen
Kreisbogen.
\end{definition}

\begin{definition}[Übergangskrümmung]
\label{def:transition-curvature}
Ein Übergangskrümmungsgesetz verbindet verschiedene Krümmungszustände über
einen ausdrücklich angegebenen Längsbereich.
\end{definition}

Seine Ableitungen machen Eigenschaften an den Anschlüssen prüfbar.
Tabelle~\ref{tab:curvature-continuity-reading} nennt nur die geometrische
Lesart; konkrete betriebliche Grenzen benötigen weiterhin anwendbares
Eisenbahnwissen, Geschwindigkeit, Überhöhung, Fahrzeugannahmen und Zweck.

\begin{table}[htbp]
\centering
\small
\caption{Stetigkeitslesarten an einer Elementgrenze $s_b$}
\label{tab:curvature-continuity-reading}
\begin{tabular}{@{}p{0.16\textwidth}p{0.25\textwidth}p{0.49\textwidth}@{}}
\toprule
Größe & Frage an der Grenze & Was die Antwort offenlegt \\
\midrule
$\kappa$ & Stimmen linker und rechter Wert überein? & Ob Krümmung stetig ist; ein Sprung führt zu einer unmittelbaren Änderung der seitlichen geometrischen Anforderung. \\
$\kappa'$ & Stimmen die Krümmungsraten überein? & Ob die Krümmungsentwicklung ohne Ratensprung ein- und ausläuft; bei Geschwindigkeit beeinflusst dies den zeitlichen Verlauf der Querreaktion. \\
$\kappa''$ & Stimmen die Ableitungen der Rate überein? & Ob die Übergangsentwicklung höherer Ordnung glatt ist; ihre Bedeutung für Fahrzeugreaktion oder Komfort bleibt von Geschwindigkeit, Überhöhung und Modell abhängig. \\
\bottomrule
\end{tabular}
\end{table}

\begin{definition}[Krümmungsstetigkeit]
\label{def:curvature-continuity}
Ein Krümmungsgesetz ist $C^0$, wenn $\kappa$ stetig ist, $C^1$, wenn
$\kappa'$ existiert und stetig ist, und $C^2$, wenn $\kappa''$ auf dem
angegebenen Bereich existiert und stetig ist.
\end{definition}

Positionsstetigkeit allein beantwortet diese Fragen nicht. Das nächste Kapitel
klassifiziert die Gesetze nach genau diesen sichtbaren Eigenschaften, bevor
das Berlinish-Kapitel sie auf den Übergangsbereich anwendet.

\section{Dünn besetzte Krümmungsrepräsentation}

\begin{definition}[Dünn besetzte Krümmungsrepräsentation]
\label{def:sparse-curvature-representation}
Eine dünn besetzte Repräsentation beschreibt Krümmung durch:
\begin{itemize}
\item feste Krümmungselemente,
\item Übergangselemente,
\item und kompakte Parametersätze.
\end{itemize}
\end{definition}

Innerhalb von AIM werden dünn besetzte Alignments als kompakte
Krümmungsprogramme interpretiert. Das dünn besetzte Modell speichert
Krümmungsstruktur, Übergangsfamilien und Rekonstruktionsregeln anstelle dichter
geometrischer Punktwolken.

Dadurch eignen sich dünn besetzte Alignments für Rekonstruktion,
Importnormalisierung, Laufzeitbewertung, Realisierungsvergleich und
Optimierung.

\section{Hybride Krümmungsrepräsentation}

\begin{definition}[Hybrides Krümmungsmodell]
\label{def:hybrid-curvature-model}
Ein hybrides Krümmungsmodell kombiniert:
\begin{itemize}
\item parametrische Krümmungsstruktur,
\item abgetastete Korrekturdaten,
\item und lokale Verfeinerung.
\end{itemize}
\end{definition}

Hybride Repräsentationen ermöglichen kompakte Speicherung, effiziente
Rekonstruktion, realisierungsbewusste Korrektur, Integration gemessener
Zustände und skalierbare Laufzeitbewertung.

Sie sind besonders relevant, wenn gemessene oder realisierte Geometrie durch
ein rein analytisches Zielmodell nicht angemessen repräsentiert werden kann.

\section{Normierter Krümmungsraum}

\begin{definition}[Normiertes Krümmungsgesetz]
\label{def:normalized-curvature-law}
Ein normiertes Krümmungsgesetz ist definiert auf
\[
w\in[0,1].
\]
\end{definition}

Normierte Repräsentationen trennen geometrische Form von physischer
Skalierung. Dies ermöglicht wiederverwendbare Übergangsfamilien,
Lookup-basierte Übergänge, Familieninterpolation und normierte Optimierung.
Physische Krümmungsgesetze entstehen aus normierten Gesetzen durch Skalierung
von Länge und Krümmungsamplitude.

\section{Krümmungsableitungen}

\begin{definition}[Krümmungsableitungen]
\label{def:curvature-derivatives}
Wichtige Größen höherer Ordnung sind:
\[
\kappa'(s),\qquad \kappa''(s),\qquad \kappa'''(s).
\]
\end{definition}

Diese Größen beeinflussen Ruck, Realisierungsverhalten, dynamische Glätte,
Fahrzeugreaktion und Übergangsqualität. Krümmungsableitungen bilden die
Brücke von intrinsischer Geometrie zur Laufzeitdynamik.

\section{Frequenzinterpretation}

Krümmungsgesetze können auch spektral interpretiert werden. Niederfrequente
Krümmungsanteile bestimmen die globale geometrische Struktur, während
hochfrequente Anteile Realisierungsempfindlichkeit, numerische Instabilität,
Messempfindlichkeit und dynamische Rauheit beeinflussen.

\begin{proposition}[Spektrales Realisierungsprinzip]
\label{prop:spectral-realization-principle}
Realisierung wirkt im Krümmungsraum näherungsweise als Tiefpassfilter.
\end{proposition}

Diese spektrale Sicht verbindet den Krümmungsraum unmittelbar mit Bau,
Instandhaltung und physischer Realisierung.

\section{Krümmungsraum und Laufzeit}

Die Laufzeitbewertung arbeitet primär auf Krümmungsstrukturen und nicht
unmittelbar auf dichter Geometrie. Projektion, Rekonstruktion,
Frame-Bewertung, pose3-Bewertung und Dynamik werden aus der Laufzeitbewertung
krümmungsbasierter Zustände abgeleitet. Der Krümmungsraum ist somit nicht nur
ein Entwurfsraum, sondern auch der rechnerische Zustandsraum der
Laufzeitrekonstruktion.

\section{Krümmungsraum und Realisierung}

Realisierung transformiert Zielkrümmung in physisch realisierte Krümmung:
\[
\kappa_{\mathrm{real}}
=
\mathcal R_{\kappa}(\kappa_{\mathrm{target}}).
\]
Der Realisierungsoperator kann den Krümmungsverlauf abschwächen, deformieren
oder lokal verändern. Der physisch relevante Entwurfsgegenstand ist folglich
nicht nur das Zielkrümmungsgesetz, sondern auch sein erwarteter realisierter
Krümmungszustand.

\section{Krümmungsraum und Optimierung}

Optimierung innerhalb von AIM wirkt primär im Krümmungsraum. Typische
Optimierungsvariablen sind Übergangsparameter, Krümmungsamplituden,
Übergangslängen, Krümmungsanker und dünn besetzte Krümmungskorrekturen.

AIM-Optimierung kann daher als Suche nach zulässigen, realisierbaren und
betrieblich geeigneten Pfaden im Krümmungsraum interpretiert werden.

\section{Kanonische Interpretation}

\begin{proposition}
Innerhalb von AIM ist Eisenbahn-Alignment-Geometrie grundlegend ein Problem
des Krümmungsraums.
\end{proposition}
\begin{proposition}
Geometrie im Positionsraum wird als rekonstruierte, aus dem
Krümmungsverlauf abgeleitete Geometrie interpretiert.
\end{proposition}
\begin{proposition}
Übergangsbögen werden als gesteuerte Übergänge im Krümmungsraum
interpretiert.
\end{proposition}
\begin{proposition}
Realisierung, Laufzeitbewertung und Optimierung wirken auf Strukturen des
Krümmungsraums, bevor sie auf visualisierte Positionsgeometrie wirken.
\end{proposition}

\section{Verbindung zu AIM}

\begin{summary}
Der Krümmungsraum bildet die intrinsische geometrische Grundlage von AIM.

In diesem Framework wirkt Krümmung als primäre Geometrie, werden dünn
besetzte Modelle zu Krümmungsprogrammen und Übergangsbögen zu
Krümmungsübergängen. Laufzeitsysteme arbeiten auf intrinsischer
Krümmungsstruktur, Realisierung transformiert Krümmungszustände und
Optimierung sucht innerhalb von Krümmungsraumstrukturen.

Diese Interpretation vereinigt Geometrie, Realisierung, Laufzeitbewertung,
Messinterpretation und Optimierung in einem intrinsischen mathematischen
Framework.
\end{summary}
